\section{Evaluation Metrics}\label{sec:method:metrics}

Inference returns, for each window, an emitter labelling $(\hat{e}_1,\ldots,\hat{e}_L)$ and a mode labelling $(\hat{m}_1,\ldots,\hat{m}_L)$ (\cref{eq:method:cluster_decoding}).
These are compared with the simulator labels $(e_n,m_n)$ of \cref{eq:method:label_pair} on the same window.
Cluster names are assigned independently per window and per branch, so only scores that are invariant to relabelling are valid.
The definitions are those of \cref{sec:bg:cluster_metrics}.

Each branch is scored separately: emitter agreement from $(e_n,\hat{e}_n)$ and mode agreement from $(m_n,\hat{m}_n)$.
Global mode labels enter training as an inductive bias that pulls same-mode embeddings together.
They are not a catalogue of names at inference, and the mode score remains a partition comparison.

The reported scores are $\ARI$, $\AMI$, homogeneity, and completeness, computed on every test window.
$\ARI$ summarises pair agreement and stays comparable when the number of active emitters or modes changes between windows.
$\AMI$ measures how much one labelling tells about the other.
It is less dominated by the largest clusters.
Pulse counts per emitter are often unbalanced, so that difference is useful here.
Homogeneity and completeness are reported separately rather than only through the V-measure.
A high homogeneity with low completeness means true emitters or modes were split; the reverse means distinct sources were merged.

Two-dimensional projections of $\{\vect{z}^{\mathrm{em}}_n\}$ and $\{\vect{z}^{\mathrm{md}}_n\}$, coloured by $(e_n,m_n)$, show whether same-label pulses sit together in each embedding.
Protocols and numbers are in \cref{ch:experiments}.
